% Body-only transcription: input this file after xepersian and macro definitions.
\sourcepage{1}{20261005\_005150919.jpg}
\section{۹ ـ سردرد و افزایش فشار داخل جمجمه}
\subsection{میگرن}
\begin{itemize}
\item معمولاً شروع \marktext{قبل از ۴۰ سالگی} / در \marktext{زنان} شایع‌تر.
\item شروع پس از ۵۰ سالگی ناشایع است.
\item در ۸۰\% موارد سابقهٔ خانوادگی مثبت وجود دارد.
\item معیارهای تشخیص: (a) حداقل \marktext{۵ حمله} سردرد مشابه که هرکدام \marktext{۴ تا ۷۲ ساعت} طول بکشد + ۲ مورد زیر: ۱. یک‌طرفه بودن؛ ۲. ضربان‌دار بودن؛ ۳. تهوع و استفراغ؛ ۴. فوتوفوبی و فونوفوبی؛ ۵. سردرد شدید.
\item انواع میگرن: \marktext{کلاسیک}: همراه با \marktext{اورا} است. \marktext{شایع}: بدون اورا / \marktext{شایع‌ترین نوع سردرد}.
\item اورای میگرن: مکانیسم \en{Spreading Depression}؛ کاهش خون‌رسانی به کورتکس لوب اکسیپیتال رخ می‌دهد.
\item اورای بینایی: \marktext{شایع‌ترین اورا} در میگرن؛ هالوسینیشن، اسکوتوم، اسکوتوم مرکزی بینایی تا همی‌آنوپی (میگرن با اورا). در ۱۰\% موارد به شکل اسکوتوم پاراسنترال که سپس شبیه \en{C} می‌شود. در برخی بیماران: به‌صورت خطوط زیگزاگ.
\item اورای حسی: بی‌حسی و مورمور شدن صورت و اندام‌ها که می‌تواند یک‌طرفه یا دوطرفه باشد.
\item طول مدت اورا: \marktext{۲۵ تا ۳۰ دقیقه} است.
\end{itemize}
\subsubsection{تفاوت اورای میگرن با \en{TIA}}
\begin{itemize}
\item اورا: \marktext{وقایع مثبت نورولوژیک} مثل نور چشمک‌زن و مورمور شدن است.
\item \en{TIA}: وقایع منفی نورولوژیک مثل نابینایی یا کرختی است.
\end{itemize}

\sourcepage{2}{20261005\_005158503.jpg}
\begin{itemize}
\item عوامل تشدیدکنندهٔ میگرن: چاقی، الکل، سیگار، مصرف بیش‌ازحد مسکن‌ها، مشکلات اندوکرین مثل هایپو و هایپرتیروئیدیسم؛ [حاشیه: پریود و استرس].
\item \marktext{تشخیص میگرن نیازی به تصویربرداری ندارد}. به‌صورت نیاز بهترین روش \marktext{\en{MRI} بدون تزریق} است.
\item در میگرن: حساسیت پوست سر / \marktext{تندرنس شریان تمپورال سطحی} دیده می‌شود.
\item بیماری‌های همراه: آپنهٔ خواب، \en{RLS}؛ درمان: \en{TCA}. [یک سرواژه پیش از \en{RLS} خط خورده و خوانا نیست.]
\item میگرن در ۹۰\% موارد یک‌طرفه و در ۴۰\% موارد دوطرفه است.
\end{itemize}
\subsubsection{درمان سردرد میگرنی}
\begin{itemize}
\item غیردارویی: اتاق تاریک و آرام، اجتناب از غذاهای محرک، ورزش سبک و منظم، مصرف آب زیاد، خواب و استراحت.
\item درمان دارویی حملهٔ حاد: استامینوفن، \en{NSAID}، متوکلوپرامید، دی‌هیدروارگوتامین؛ تریپتان‌ها: سوماتریپتان، ریزاتریپتان، زولمی‌تریپتان.
\item خفیف تا متوسط: استامینوفن، \en{NSAID}.
\item متوسط تا شدید: ارگوت و تریپتان‌ها.
\end{itemize}
\subsection{سردرد ناشی از مصرف بیش‌ازحد مسکن}
\begin{latin}\noindent Medication Overuse headache\end{latin}
\begin{itemize}
\item استفاده \marktext{بیش از ۲ تا ۳ بار در هفته} از داروهای حملهٔ حاد: سردرد مزمن وابسته به دارو.
\item علی‌رغم مصرف دارو بیمار در بیشتر روزها سردرد دارد.
\item \marktext{درمان: قطع داروها}.
\end{itemize}

\sourcepage{3}{20261005\_005212674.jpg}
\subsubsection{اندیکاسیون‌های درمان پروفیلاکتیک میگرن}
\begin{itemize}
\item \marktext{حملات بیشتر از ۳ تا ۴ بار در هفته}.
\item اختلال در عملکرد روزانهٔ بیمار.
\item عدم پاسخ به داروهای ضد میگرن و مسکن.
\item کاهش کیفیت زندگی.
\end{itemize}
\subsubsection{درمان پروفیلاکتیک}
\begin{itemize}
\item به‌مدت \marktext{۶ ماه تا ۱ سال} بوده / بیمار همچنان دارو مصرف می‌کند.
\item بتابلاکرها: \marktext{پروپرانولول شایع‌ترین و مؤثرترین}.
\item \en{TCAs}: آمی‌تریپتیلین، نورتریپتیلین.
\item ضدتشنج: سدیم والپروات، توپیرامات.
\item \en{CCB}: وراپامیل، فلوناریزین.
\end{itemize}
\subsection{میگرن بازیلار}
\begin{itemize}
\item اورا به‌علت اختلال در عملکرد \marktext{ساقهٔ مغز} ایجاد می‌شود.
\item سرگیجه، دیزآرتری، دوبینی، پارستزی.
\item \marktext{۳۰ دقیقه} طول می‌کشد و سپس \marktext{سردرد ضربان‌دار} می‌شود.
\end{itemize}
\subsection{میگرن همی‌پلژیک}
\begin{itemize}
\item \marktext{همیشه در کودکی آغاز می‌شود}.
\item همی‌پارزی به‌مدت \en{20--30 min} و متعاقب آن سردرد.
\item سمت درگیر ممکن است تغییر کند.
\item در نوع شدید، پس از سردرد نیز همی‌پلژی به‌مدت چند روز تا چند هفته می‌ماند.
\item \marktext{در میگرن بازیلار و همی‌پلژیک استفاده از تریپتان‌ها ممنوع است}.
\end{itemize}
\subsection{میگرن افتالموپلژیک}
\begin{itemize}
\item همیشه در کودکی آغاز می‌شود.
\item سردرد همراه با دوبینی، پتوز، و استفراغ.
\item ممکن است ۴ تا ۱۰ روز طول بکشد.
\item میگرن سردرد \marktext{شدید همان سمت} و طی چند ساعت \marktext{فلج کامل عصب ۳}، میدریاز و عدم پاسخ به نور.
\end{itemize}

\sourcepage{4}{20261005\_005219492.jpg}
\begin{itemize}
\item \marktext{فلج چشمی از چند روز تا ۲ ماه}.
\item پس از چند حمله ممکن است علائم پارزی عصب ۳ باقی بماند.
\end{itemize}
\subsubsection{میگرن در بارداری}
\begin{itemize}
\item میگرن در بارداری: $1/3$ تشدید، $1/3$ تخفیف، $1/3$ را تغییری نمی‌کند.
\item طی سه‌ماههٔ دوم و سوم، تعداد حملات کاهش می‌یابد.
\item احتمال بهبود میگرن در افرادی که سردرد آن‌ها با تغییرات هورمون مرتبط است بیشتر می‌باشد.
\item مثل کسانی که با شروع اولین قاعدگی سردرد را تجربه کرده‌اند، و یا افرادی که میگرن قاعدگی دارند.
\item اگر اولین سردرد در حاملگی رخ دهد یا شدت پیدا کند: واسکولیت، تودهٔ مغزی، به‌ویژه مننژیوم، آنوریسم، \en{AVM}، ادم هیپوفیز.
\item ترومبوز وریدهای مغزی در بارداری شایع‌تر است.
\end{itemize}
\subsubsection{درمان میگرن در بارداری}
\begin{itemize}
\item سردرد حاد: مسکن معمولی؛ استامینوفن ارجح است.
\item اگر تهوع و استفراغ: متوکلوپرامید، کلرپرومازین.
\item ارگوتامین و تریپتان‌ها در بارداری کنتراندیکه هستند.
\item اگر سردرد خیلی شدید باشد: اپیوئید می‌دهیم.
\end{itemize}
\subsubsection{پروفیلاکسی میگرن در حاملگی}
\begin{itemize}
\item کنترل منظم روزانه: خواب منظم، اجتناب از غذاهای محرک، ورزش.
\item بتابلاکر، \en{TCA}، ضدتشنج در بارداری ممنوع است.
\item بتابلاکر به‌ویژه در سه‌ماههٔ آخر: \en{IUGR}.
\end{itemize}

\sourcepage{5}{20261005\_005230790.jpg}
\subsection{سردرد تنشی (\en{Tension})}
\begin{itemize}
\item سردرد مزمن / \marktext{شایع‌ترین نوع سردرد}.
\item علائم فوکال نورولوژیک و علائم گوارشی میگرن وجود ندارد.
\item در زنان و در میان‌سالی شایع‌تر است / ولی در هر سنی دیده می‌شود.
\item آستانهٔ درد در این افراد پایین‌تر است.
\item متعاقب فشار روحی.
\item تظاهرات بالینی: دوطرفه در نواحی پس‌سری + گردن، تمپورال، فرونتال.
\item به‌شکل احساس فشار، درد مبهم، سفتی نواری‌شکل، وجود باند دور سر.
\item از \en{30 min} تا چند روز.
\item معمولاً فرد می‌تواند کار کند و زیاد شود ولی باعث اختلال خواب نمی‌گردد.
\item با تهوع، استفراغ، فتوفوبی و فونوفوبی همراهی ندارد.
\item با فعالیت فیزیکی مرتبط نیست.
\item ارتباط نزدیکی با افسردگی و اضطراب دارد.
\item در موارد طول‌کشیده می‌تواند مثل میگرن یک‌طرفه و ضربان‌دار گردد.
\item بیماری‌های همراه: آپنهٔ خواب، \en{RLS}؛ درمان: \en{TCA}.
\item \en{DDX}: آرتریت تمپورال؛ اگر در بیمار مسن با سردرد تنشی باشد \en{ESR} چک می‌شود.
\item ضایعات اینتراکرانیال: در هر بیمار \unclear{علامت مقایسه کنار ۵۰ سال؛ جهت مقایسه در دست‌نویس مبهم است}: \en{CT} و \en{MRI} (اطمینان دادن).
\item آتیپیک میگرن: معاینهٔ نرمال و رادیولوژی نرمال افتراق.
\end{itemize}
\subsubsection{درمان}
\begin{itemize}
\item خفیف و گاهگاهی: استامینوفن، \en{NSAID}، استراحت، ورزش.
\item موارد مزمن: مؤثرترین داروی پیشگیری: آمی‌تریپتیلین؛ کاهش حساسیت به درد بیمار، هنگام خواب \en{50--150 mg/day}.
\item اگر آمی‌تریپتیلین کنتراندیکه باشد از سایر \en{TCAs}، گاباپنتین، توپیرامات استفاده می‌کنیم.
\item در درمان \en{Tension}: پروپرانولول و ارگوتامین مؤثر نیست.
\item همیشه بهتر است ورزش، یوگا و ریلکسیشن بدهیم.
\end{itemize}

\sourcepage{6}{20261005\_005236315.jpg}
\subsection{سردرد خوشه‌ای (کلاستر)}
\begin{itemize}
\item دردناک‌ترین سردرد راجعه است.
\item در مردان ۳ تا ۹ برابر شایع‌تر است / معمولاً در ۲۰ تا ۴۰ سالگی آغاز می‌شود.
\item در برخی اقوام توارث دیده شده است.
\item علائم بالینی: یک‌طرفه، بدون احساس ضربان (\en{Non-throbbing}).
\item همیشه همان سمت.
\item معمولاً در شب و ۱ تا ۲ ساعت بعد از به‌خواب‌رفتن شروع می‌شود.
\item حملات تا ۸ بار در روز تکرار می‌شود.
\item در ظرف ۵ تا ۱۰ دقیقه به \en{pick} می‌رسد و طی مدت ۴۵ دقیقه تا ۲ ساعت ادامه می‌یابد.
\item محل درد رترواربیتال یا تمپورال است.
\item انگار مته وارد چشم بیمار شده.
\item علائم اتونوم هم دیده می‌شود: اشک‌ریزش، احتقان ملتحمه، گرفتگی بینی، پتوز، میوز (سندرم هورنر) در همان سمت.
\item بیماری‌های همراه: زخم دئودنوم، افزایش اسید معده (سندرم زولینگر الیسون).
\item الکل موجب شعله‌ور شدن حملات می‌شود.
\end{itemize}
\subsubsection{افتراق کلاستر از میگرن}
\begin{itemize}
\item تهوع، فتوفوبی، فونوفوبی در میگرن شایع‌تر است.
\item استفراغ، اورا در کلاستر وجود ندارد.
\item در کلاستر فرد تمایلی به درازکشیدن ندارد و قدم‌زدن یا نشستن را ترجیح می‌دهد.
\item پریودیک بودن ویژگی اصلی کلاستر است.
\item هر سال ۱ تا ۲ پریود؛ هر پریود ۴ تا ۹ هفته؛ ۱ تا ۳ بار در روز.
\item زمان بروز حملات در روزهای مختلف یکسان است.
\item طول حملات در هر پریود متفاوت و مدت هر پریود نیز ثابت است.
\end{itemize}

\sourcepage{7}{20261005\_005241413.jpg}
\subsubsection{درمان کلاستر}
\begin{itemize}
\item حملهٔ حاد: اکسیژن ۱۰۰\% به‌مدت ۱۰ تا ۱۵ دقیقه.
\item سوماتریپتان \en{SC} / دی‌هیدروارگوتامین \en{SC} یا \en{IN}.
\item پروفیلاکسی: داروهایی که به‌سرعت اثر می‌کنند.
\item پردنیزولون ۲۴ تا ۲۸ روز.
\item ارگوتامین \en{2 mg} \unclear{عدد دوز: ۲ یا ۳} شب‌ها / دی‌هیدروارگوتامین.
\item داروهای بالا نیاز به زمان دارند.
\item وراپامیل / لیتیوم / متی‌سرژید / ایندومتاسین.
\item ماساژ یا یخ گاهی باعث بهبود کلاستر می‌شود.
\end{itemize}
\subsection{نورالژی تری‌ژمینال}
\begin{itemize}
\item درد حمله‌ای در مسیر یک یا چند شاخهٔ عصب تری‌ژمینال.
\item در زنان شایع‌تر / ۹۰\% بعد از ۴۰ سالگی / با افزایش سن احتمال ایجاد بیشتر.
\item اتیولوژی: ایدیوپاتیک؛ تحت فشار قرار گرفتن ریشهٔ خلفی عصب توسط یک حلقهٔ عروقی / معاینهٔ حسی ـ حرکتی عصب \en{NL} است.
\item سایر علل: \en{MS}، آنوریسم شریان بازیلار، تومورهای \en{CPA} مانند شوانوم و مننژیوم.
\item وجود هر اختلالی در حس و حرکت عصب نشان‌دهندهٔ وجود ضایعه است.
\item اگر دوطرفه باشد احتمال \en{MS} خیلی بالاست.
\item علائم بالینی: درد حمله‌ای / چند ثانیه تا ۲ دقیقه / بیشتر شاخهٔ ۲ و ۳ درگیر می‌شود / اغلب یک‌طرفه است.
\item درد به‌صورت خنجری، شوک الکتریکی، تیرکشنده.
\item در روز و شب بارها تکرار می‌شود ولی موقع خواب بروز نمی‌ـ [ادامه در صفحهٔ بعد].
\end{itemize}

\sourcepage{8}{20261005\_005246762.jpg}
\begin{itemize}
\item [ادامهٔ جملهٔ صفحهٔ قبل:] می‌شود.
\item در صورت تکرار حملات بیمار دچار کاهش وزن، دهیدراتاسیون و افسردگی می‌شود.
\item عوامل \en{trigger}: درد با تحریک پوست، مخاط دهان و زبان شروع می‌شود؛ جویدن، مسواک زدن، حرف زدن، تراشیدن ریش، خمیازه.
\end{itemize}
\subsubsection{درمان نورالژی تری‌ژمینال}
\begin{itemize}
\item خط اول کاربامازپین با دوز \en{900--1200 mg} در روز.
\item فنی‌توئین، باکلوفن، کلونازپام، گاباپنتین.
\item موارد مقاوم: جراحی.
\item در موارد زیر نورالژی تری‌ژمینال با شک به \en{MS} برای بیمار \en{MRI} درخواست می‌کنیم: بیمار جوان باشد؛ نورالژی دوطرفه باشد؛ هرگونه اختلال حسی یا حرکتی در عصب تری‌ژمینال.
\end{itemize}
\subsection{آرتریت تمپورال}
\begin{itemize}
\item واسکولیت شریان‌های متوسط تا بزرگ / در افراد مسن (بالای ۵۰ سال). [حاشیه: انسداد و تنگی عروق].
\item اگر سریع تشخیص داده نشود: منجر به کوری می‌شود.
\item در زنان شایع‌تر / اکثر بیماران بالای ۶۵ سال.
\item علائم بالینی: سردرد (شایع‌ترین)، ضربان‌دار یا غیرضربانی، معمولاً یک‌طرفه؛ لوکالیزه در شریان تمپورال سطحی یا اکسیپیتال؛ در شب‌ها بدتر می‌شود.
\item پلی‌میالژیا روماتیکا در بیش از ۵۰\% بیماران: درد شانه و محل پروگزیمال در عضلات، خستگی، بی‌حالی، ضعف، تب، \en{Malaise}.
\item نشانه‌های التهاب شریان تمپورال سطحی (۵۰\%): اریتم، تندرنس، ندولاریته، کاهش نبض شریان، افزایش ضخامت.
\item \en{Jaw claudication}: به‌علت ایسکمی عضلات جونده، درد با سفتی عضلات در ۴۰\% بیماران. [حاشیه: هنگام غذا خوردن].
\item کوری گذرا: شایع‌ترین عارضهٔ آرتریت تمپورال؛ ۵۰\% کوری دائمی. [حاشیه: \en{Bl}].
\item نوروپاتی ایسکمیک عصب اپتیک.
\end{itemize}

\sourcepage{9}{20261005\_005253981.jpg}
\begin{itemize}
\item فلج حرکات چشم: در ۱۰\% بیماران به‌دلیل فلج عصب ۳ یا ۶.
\item سکتهٔ مغزی در این بیماران نادر است.
\item سایر علائم: نوروپاتی، میلوپاتی، \en{stroke} مغزی، دوبینی، دیسفاژی، کلودیکیشن.
\item یافته‌های آزمایشگاهی: \en{ESR} معمولاً بالای ۵۰ تا ۷۵ است / در ۳\% موارد طبیعی است.
\item افزایش \en{CRP} نسبت به \en{ESR} ملاک حساس‌تری است.
\item آنمی، ترومبوسیتوز خفیف، لکوسیتوز با ارجحیت \en{PMN}.
\item در بیمار سالخورده با سردرد مداوم: \en{ESR} را چک کنیم.
\item تشخیص قطعی با بیوپسی از شریان تمپورال سطحی است.
\item به محض شک به آرتریت تمپورال: درمان با دوز بالای پردنیزولون \en{40--60 mg/daily}.
\item بهتر است بیوپسی قبل از شروع درمان انجام شود.
\item شروع درمان سریع را نباید فدای انجام بیوپسی نمود.
\item بررسی پاسخ به درمان: علائم بالینی و \en{ESR} کنترل می‌شود.
\item در چند روز اول درمان سردرد بهبود می‌یابد و \en{ESR} طبیعی می‌شود!
\item طول مدت درمان با کورتون دوز بالا حداقل ۱ ماه است.
\end{itemize}
\subsection{سردرد رعدآسا (\en{Thunderclap})}
\begin{itemize}
\item سردرد شدید و ناگهانی بدون هیچ‌گونه علائم هشدار قبلی.
\item ظرف یک دقیقه به \en{pick} می‌رسد.
\item اتیولوژی: خون‌ریزی ساب‌آراکنوئید (\en{SAH})، ترومبوز سینوس وریدی، آپوپلکسی هیپوفیز، دایسکشن شریان‌های گردن، کریز حاد فشار خون، مننژیت، افت \en{BP} اینتراکرانیال، در اثر نشت \en{CSF}؛ سندرم انقباض برگشت‌پذیر عروق مغزی؛ انفارکتوس مغزی ناشی از آمبولی.
\end{itemize}

\sourcepage{10}{20261005\_005300572.jpg}
\begin{itemize}
\item علائم \en{SAH}: سردرد ناگهانی / درد گردن / کاهش هوشیاری / خون‌ریزی سطح ملتحمه (ساب‌هیالوئید).
\item هر بیمار با سردرد شدید و ناگهانی: \en{CT} اورژانسی.
\item اگر \en{CT} طبیعی بود اقدام بعدی: \en{LP}.
\end{itemize}
\subsection{سردرد ناشی از ضایعات فضاگیر مغزی}
\begin{itemize}
\item تومور، آبسه، ضایعات گرانولوماتوز: افزایش \en{ICP}؛ سردرد غیراختصاصی.
\item ۵۰ تا ۷۵\% مبتلایان به تومور مغزی با سردرد تظاهر پیدا می‌کنند.
\item در ضایعات بدخیم سریع و ضایعات اینفراتنتوریال زودتر بروز می‌کند.
\item افزایش \en{ICP} منجر به سردرد ژنرالیزه می‌شود که در ناحیهٔ اکسیپیتال و گردن شدیدتر است.
\item ضایعات سوپراتنتوریال: سردرد ناحیهٔ فرونتال یا تمپورال.
\item ضایعات اینفراتنتوریال: سردرد ناحیهٔ اکسیپیتال و گردن.
\item ویژگی سردرد ناشی از ضایعات فضاگیر: با مانور والسالوا تشدید می‌شود؛ معمولاً با تهوع و استفراغ همراه هستند؛ پس از بیدار شدن از خواب در صبح شدیدتر است.
\end{itemize}
\subsubsection{علائم هشداردهنده}
\begin{itemize}
\item سردرد تحت‌حاد و پیشرونده؛ معاینهٔ نورولوژیک غیرطبیعی.
\item شروع سردرد بعد از ۴۰ سالگی؛ همراهی با تشنج.
\item تغییر الگوی سردرد؛ همواره یک‌طرفه باشد.
\item سردردی که بیمار را از خواب بیدار کند؛ عدم پاسخ به مسکن معمولی.
\item همراهی با تهوع و استفراغ، ولی شواهد بیماری سیستمیک نباشد.
\item تشدید با مانور والسالوا و تغییر وضعیت؛ سردرد رعدآسا.
\item تشدید در صبح هنگام بیدار شدن از خواب.
\item هیدروسفالی، اختلال هوشیاری.
\end{itemize}

\sourcepage{11}{20261005\_005306249.jpg}
\begin{itemize}
\item ایجاد سردرد جدید در بیماران سرطان و نقص ایمنی.
\item همراهی با اختلالات اندوکرین (مثل گالاکتوره).
\item در صورت وجود علائم هشداردهنده: \en{MRI}.
\end{itemize}
\subsection{هایپرتنشن ایدیوپاتیک داخل جمجمه}
\begin{itemize}
\item سودوتومور سربری.
\item افزایش \en{ICP}، سردرد، ادم پاپی.
\item \en{CT}، \en{MRI}، آنالیز \en{CSF} طبیعی است.
\item احتمالاً به‌دلیل انسداد عملکردی در جریان خروجی خون از سینوس‌های وریدی.
\item در زنان جوان شایع‌تر است / در کودکان و نوجوانان نیز دیده می‌شود.
\item \en{RF}: چاقی، اختلالات قاعدگی، داروها.
\item هیپوکلسمی، هایپرویتامینوز \en{A} (ایزوترتینوئین)، نالیدیکسیک اسید، تتراسایکلین، وریکونازول، سایمتیدین، پردنیزولون، تاموکسیفن، نیتروفورانتوئین، سولفونامیدها، کینولون‌ها.
\end{itemize}
\subsubsection{علائم بالینی}
\begin{itemize}
\item سردرد: معمولاً اکسیپیتال، ژنرالیزه / هنگام صبح شدیدتر / با مانور والسالوا بدتر.
\item ادم پاپی، دوبینی (فلج عصب ۶)، احساس گیجی و منگی، وزوز گوش.
\item تاری دید موقتی: در هنگام تغییر وضعیت از نشسته به ایستاده.
\item محدود شدن میدان بینایی و بزرگ شدن نقطهٔ کور؛ در طولانی‌مدت کوری دائمی.
\end{itemize}
\subsubsection{تشخیص}
\begin{itemize}
\item \en{MRI} یا \en{CT}: بطن‌ها کوچک یا شکاف‌مانند (\en{Slitlike}).
\item وقتی وجود توده \en{R/O} شد: \en{LP}؛ فشار \en{CSF} $>25$ سانتی‌متر آب \unclear{علامت دست‌نویس به‌شکل ۲۵ < فشار نوشته شده}.
\item قند \en{NL}، \en{PMN} ندارد، \en{Pro}، \en{NL} یا کاهش.
\end{itemize}
\subsubsection{درمان}
\begin{itemize}
\item استازولاماید، پردنیزولون.
\item با مهار کربنیک انهیدراز تولید \en{CSF} را پایین می‌کند.
\item \en{LP} مکرر / جراحی شانت: لومبوپریتونئال؛ کاهش وزن؛ سوراخ کردن غلاف عصب اپتیک.
\end{itemize}
